% !TEX root = ../main.tex
\chapter{How Matter Shapes Spacetime}\label{ch:matter-geometry}

\begin{goals}
\begin{itemize}
\item Describe the energy, momentum and stress of matter with the
  stress-energy tensor, and say what each of its components means.
\item Compute the energy density that any observer measures.
\item Write down the stress-energy of dust, fluids, radiation, vacuum energy
  and tension, and explain why tension counts as negative pressure.
\item State Einstein's equation, convert between geometric and everyday units,
  and recover Newton's law of gravity from it.
\item Carry out the spacetime engineer's central calculation: from a chosen
  metric to the matter it demands.
\item Recognize the first signs of an exotic demand.
\end{itemize}
\end{goals}

\Cref{ch:geometry} gave us the left side of Einstein's equation, the geometry.
This chapter supplies the right side, the description of matter, and then puts
the two together. By the end you will be able to take any metric, compute the
matter it requires, and read the result the way a spacetime engineer does: as
a list of demands that some physical material would have to meet.

\section{Energy depends on who measures it}\label{sec:energy-observers}

In Newtonian physics, the mass of a brick is a fixed number. In relativity the
brick's \emph{energy} depends on the observer: someone moving past it sees it
carry kinetic energy as well as its rest energy. To describe matter in a way
all observers can use, we need to say how the measured energy changes with the
observer.

A particle of mass $m$ with four-velocity $u^{\mu}$ has \emph{four-momentum}
$p^{\mu} = m\,u^{\mu}$. An observer with four-velocity $w^{\mu}$ measures the
particle's energy to be
\begin{equation}
  E = -g_{\mu\nu}\,p^{\mu}w^{\nu} \equiv -p_{\mu}w^{\mu}.
  \label{eq:measured-energy}
\end{equation}
In the second form we have \emph{lowered an index}: $p_{\mu} = g_{\mu\nu}p^{\nu}$.
Lowering indices with the metric is the bookkeeping that turns a pair of
vectors into a number every observer agrees on. For an observer at rest in flat
spacetime, $w^{\mu} = (1,0,0,0)$, and \cref{eq:measured-energy} gives $E = p^{0}$,
the time component of the four-momentum. For a particle at rest relative to the
observer, $E = m$: its rest energy, $mc^{2}$ in ordinary units.

\section{The stress-energy tensor}\label{sec:stress-energy}

A single particle is described by one four-momentum. A continuous material,
such as a gas, a fluid or a stretched rubber band, needs more. At each event
we want to know how much energy is present per unit volume, how fast energy is
flowing, how much momentum there is, and what forces neighbouring parts of the
material exert on each other. All of this fits into one symmetric $4\times4$
array, the \emph{stress-energy tensor} $T^{\mu\nu}$.

\begin{definition}[Stress-energy tensor]
In the frame of an observer, with coordinates $t, x, y, z$,
\begin{itemize}
\item $T^{tt}$ is the \emph{energy density}, the energy per unit volume;
\item $T^{ti} = T^{it}$ is the \emph{energy flux} in direction $i$, which equals
  the density of momentum in direction $i$;
\item $T^{ij}$ is the \emph{stress}: the flux of $i$-momentum across a surface
  facing the $j$ direction. Its diagonal entries are \emph{pressures}, and its
  off-diagonal entries are \emph{shear stresses}.
\end{itemize}
\end{definition}

\begin{figure}[htbp]
\centering
\begin{tikzpicture}[scale=0.9, every node/.style={font=\small}]
  \fill[keygold!35] (0,3) rectangle (1,4);
  \fill[examplegreen!20] (1,3) rectangle (4,4);
  \fill[examplegreen!20] (0,0) rectangle (1,3);
  \fill[goalblue!12] (1,0) rectangle (4,3);
  \foreach \i in {1,2,3} { \fill[goalblue!30] (\i,3-\i) rectangle (\i+1,4-\i); }
  \draw[step=1, gray] (0,0) grid (4,4);
  \draw[thick] (0,0) rectangle (4,4);
  \foreach \l [count=\i from 0] in {t,x,y,z} {
    \node at (\i+0.5,4.3) {$\l$};
    \node at (-0.3,3.5-\i) {$\l$};
  }
  \node[align=left, anchor=west] at (4.6,3.5) {energy density $T^{tt}$};
  \node[align=left, anchor=west, text=examplegreen!80!black] at (4.6,2.6) {energy flux = momentum density\\$T^{ti} = T^{it}$};
  \node[align=left, anchor=west, text=goalblue] at (4.6,1.3) {stress $T^{ij}$: pressures on the\\diagonal, shear off the diagonal};
\end{tikzpicture}
\caption{The components of the stress-energy tensor in an observer's frame. The
array is symmetric, so ten numbers describe the matter at each event.}
\label{fig:stress-energy}
\end{figure}

Different observers split the same tensor into energy, flux and stress
differently, just as observers in relative motion split an interval into time
and distance differently. The rule that makes this precise generalizes
\cref{eq:measured-energy}: an observer with four-velocity $w^{\mu}$ measures the
energy density
\begin{equation}
  \rho_{w} = T_{\mu\nu}\,w^{\mu}w^{\nu},
  \label{eq:observed-density}
\end{equation}
where $T_{\mu\nu} = g_{\mu\alpha}g_{\nu\beta}T^{\alpha\beta}$ is the tensor with
both indices lowered. This formula is one of the most used in the book. When a
geometry demands matter, the first question is always what energy density
each observer would measure.

\begin{example}{A moving cloud of dust}\label{ex:dust}
A cloud of dust has particles of mass $m$ at rest relative to one another, with
$n$ particles per unit volume in their own rest frame. Find the components of
$T^{\mu\nu}$ measured by an observer past whom the cloud moves at speed $v$
along $x$.

\solution In the cloud's rest frame each particle has energy $m$, so the energy
density is $\rho_{0} = nm$, and there is no flux and no stress. The moving
observer sees two changes. Each particle has energy $\gamma m$, with
$\gamma = 1/\sqrt{1-v^{2}}$. The cloud is also length-contracted along $x$, so
the number density rises to $\gamma n$. The energy density is therefore
\begin{equation*}
  T^{tt} = \gamma^{2}\rho_{0}.
\end{equation*}
Energy flows along $x$ at speed $v$, so the energy flux is
$T^{tx} = \gamma^{2}\rho_{0}\,v$. Momentum in the $x$ direction,
$\gamma^{2}\rho_{0}v$ per unit volume, is carried along $x$ at speed $v$, so the
momentum flux is $T^{xx} = \gamma^{2}\rho_{0}\,v^{2}$. All other components
vanish. These are exactly the components of
\begin{equation}
  T^{\mu\nu} = \rho_{0}\,u^{\mu}u^{\nu}
  \label{eq:dust}
\end{equation}
with $u^{\mu} = \gamma(1, v, 0, 0)$.
\discussion The energy density of the same dust is $\rho_{0}$ for one observer
and $\gamma^{2}\rho_{0}$ for another, and the formula $T^{\mu\nu} = \rho_{0}u^{\mu}u^{\nu}$
captures both. You can check \cref{eq:observed-density} on it: an observer
moving with the dust has $w^{\mu} = u^{\mu}$, and
$T_{\mu\nu}u^{\mu}u^{\nu} = \rho_{0}\,(u_{\mu}u^{\mu})^{2} = \rho_{0}$.
\end{example}

\section{A catalogue of matter}\label{sec:catalogue}

The stress-energy tensors of a few kinds of matter recur throughout this book.
Each is written below in its own rest frame, where it is simplest.

\paragraph{Perfect fluid.} A fluid with energy density $\rho$ and the same
pressure $p$ in every direction has
\begin{equation}
  T^{\mu\nu} = (\rho + p)\,u^{\mu}u^{\nu} + p\,g^{\mu\nu},
  \qquad
  T^{\mu\nu}_{\text{rest}} = \mathrm{diag}(\rho, p, p, p).
  \label{eq:perfect-fluid}
\end{equation}
Dust is a perfect fluid with $p = 0$. Gases, liquids and the interiors of stars
are perfect fluids to a good approximation.

\paragraph{Radiation.} A gas of light moving in all directions is a perfect
fluid with $p = \rho/3$.

\paragraph{Tension.} Stretch a rubber band and each part pulls on its
neighbours along the band. A pull is a negative pressure, so tension appears in
$T^{\mu\nu}$ as a \emph{negative} diagonal stress. A material under tension
$\tau$ along $z$ has $T^{zz} = -\tau$. Tension is one of the spacetime engineer's
most important ingredients, and you will see it again and again in the demands
of engineered geometries.

\paragraph{Vacuum energy.} Empty space can carry an energy density of its own,
the kind cosmologists invoke to explain the accelerating expansion of the
universe. Its stress-energy is
\begin{equation}
  T_{\mu\nu} = -\rho_{\text{vac}}\,g_{\mu\nu},
  \label{eq:vacuum-energy}
\end{equation}
which in any frame has energy density $\rho_{\text{vac}}$ and pressure
$p = -\rho_{\text{vac}}$ in every direction. Every observer measures the same
energy density: vacuum energy looks the same to everyone.

\begin{deeper}{the electromagnetic field}
In Heaviside--Lorentz units with $c = 1$, a magnetic field of strength $B$
along $z$ has energy density $B^{2}/2$, a pressure $+B^{2}/2$ across the field
lines ($T^{xx} = T^{yy}$), and a tension $B^{2}/2$ along them
($T^{zz} = -B^{2}/2$). Field lines behave like stretched elastic bands pushing
sideways on one another. The tension along the field equals the energy density:
this is the largest tension allowed by the dominant energy condition, which you
will meet in Chapter~5, and magnetic fields reach it exactly.
\end{deeper}

\begin{checkpoint}
For each of dust, radiation, a perfect fluid with $p > 0$, and vacuum energy,
compute $\rho + p$. Which of them has $\rho + p = 0$? (You will see in Chapter~5
that the sign of $\rho + p$ decides whether matter focuses light.)
\end{checkpoint}

\section{Conservation}\label{sec:conservation}

Energy and momentum are conserved: whatever flows out of a small box must be
accounted for by a loss inside it. In flat spacetime this is the statement
\begin{equation}
  \partial_{\mu}T^{\mu\nu} = 0 .
  \label{eq:conservation-flat}
\end{equation}
Its $\nu = t$ component is the continuity equation for energy,
$\partial_{t}T^{tt} + \partial_{i}T^{it} = 0$; its three spatial components say
that momentum changes only through the stresses acting on a region. In curved
spacetime the partial derivative is replaced by the \emph{covariant} derivative
$\nabla_{\mu}$, which adds Christoffel-symbol terms to account for the geometry:
$\nabla_{\mu}T^{\mu\nu} = 0$. Physically, conservation in curved spacetime says
that matter exchanges energy and momentum with the gravitational field in a
definite way, and the extra terms describe that exchange.

\section{Einstein's equation}\label{sec:einstein}

We can now write Einstein's equation precisely:
\begin{equation}
  G_{\mu\nu} = \frac{8\pi G}{c^{4}}\,T_{\mu\nu},
  \label{eq:einstein-si}
\end{equation}
with $G_{\mu\nu}$ the Einstein tensor of \cref{eq:einstein-tensor}. There are ten
equations here, one for each independent component of the two symmetric
tensors.

\begin{deeper}{why the Einstein tensor?}
Whatever sits on the left side must be consistent with conservation of energy
and momentum on the right side. The Einstein tensor has exactly this property:
its covariant divergence vanishes identically, $\nabla^{\mu}G_{\mu\nu} = 0$ for
\emph{every} metric, a result called the contracted Bianchi identity. Einstein's
equation therefore implies $\nabla^{\mu}T_{\mu\nu} = 0$ automatically. The
Einstein tensor is essentially the only combination of the metric and its first
two derivatives with this property, apart from a term proportional to the metric
itself, which acts as vacuum energy.
\end{deeper}

\subsection{Geometric units}

From now on we set $G = 1$ as well as $c = 1$. Einstein's equation then reads
\begin{equation}
  G_{\mu\nu} = 8\pi\,T_{\mu\nu}.
  \label{eq:einstein}
\end{equation}
In these \emph{geometric units} everything is measured in powers of length.
A mass $M$ becomes the length $GM/c^{2}$: the Sun is \SI{1.48}{\kilo\metre}
and the Earth \SI{4.4}{\milli\metre}. Curvature has units of one over length
squared, and so, by \cref{eq:einstein}, do energy density and stress. To
convert back, multiply a geometric energy density or stress, with lengths in
metres, by
\begin{equation}
  \frac{c^{4}}{G} = \SI{1.21e44}{\newton},
  \label{eq:conversion}
\end{equation}
to get joules per cubic metre, which are the same as pascals. Divide by $c^{2}$
to express an energy density as a mass density. This single conversion factor
is the \enquote{stiffness of spacetime} of \cref{eq:stiffness}.

\subsection{Newton's law of gravity, recovered}

\begin{example}{The Newtonian limit}\label{ex:newtonian}
Show that for the weak-field metric \cref{eq:weak-field} with a static potential
$\Phi$, the $tt$ component of Einstein's equation reduces to Poisson's equation
of Newtonian gravity, $\nabla^{2}\Phi = 4\pi\rho$.

\solution Compute the Einstein tensor of \cref{eq:weak-field} keeping only
terms of first order in $\Phi$. The Christoffel symbols are first order, so
their products in \cref{eq:riemann} are second order and can be dropped. The
calculation (\cref{prob:weak-field}) gives
\begin{equation*}
  G_{tt} \approx 2\,\nabla^{2}\Phi .
\end{equation*}
For slowly moving matter, the dominant component of the stress-energy is the
energy density, $T_{tt} \approx \rho$. Einstein's equation $G_{tt} = 8\pi T_{tt}$
then gives
\begin{equation*}
  \nabla^{2}\Phi = 4\pi\rho ,
\end{equation*}
which is Poisson's equation in units with $G = 1$.
\discussion Newtonian gravity is the weak-field, slow-motion limit of
Einstein's equation. The factor $8\pi$ in \cref{eq:einstein} is fixed by this
requirement.
\end{example}

\subsection{Pressure gravitates}

Einstein's equation says more than Newton's. For a perfect fluid at rest, a
short calculation from \cref{eq:einstein} and \cref{eq:perfect-fluid}
(\cref{prob:rho-3p}) shows that the quantity playing the role of $\rho$ in
Poisson's equation is
\begin{equation}
  \rho + 3p .
  \label{eq:active-mass}
\end{equation}
Pressure is a source of gravity alongside energy. For ordinary matter the
pressure is tiny compared with the energy density, which includes the rest
energy $mc^{2}$, so Newton never noticed. Tension, however, \emph{weakens}
gravity, and enough of it reverses the sign. Vacuum energy has
$\rho + 3p = -2\rho_{\text{vac}}$: it repels. This is why a positive vacuum
energy makes the expansion of the universe accelerate, and it is a first hint
of what engineered spacetimes can do with tension.

\begin{keyidea}
In general relativity every component of the stress-energy tensor bends
spacetime: energy, momentum and stress alike. Pressure attracts, tension
repels, and a spacetime engineer designs with all of them.
\end{keyidea}

\section{Reading the equation backwards}\label{sec:backwards}

Here is the spacetime engineer's central calculation, carried out in full. The
recipe has five steps.
\begin{enumerate}
\item Write down the metric $g_{\mu\nu}$ of the geometry you want.
\item Compute the Christoffel symbols, \cref{eq:christoffel}.
\item Compute the Riemann tensor \cref{eq:riemann}, then the Ricci tensor, the
  Ricci scalar and the Einstein tensor \cref{eq:einstein-tensor}.
\item Read off the demanded stress-energy, $T_{\mu\nu} = G_{\mu\nu}/8\pi$.
\item Interpret it: compute the energy density and stresses measured by the
  observers you care about, using \cref{eq:observed-density}.
\end{enumerate}
For all but the simplest metrics, steps 2 and 3 are best done with computer
algebra. The examples below are simple enough to follow by hand, and each
teaches something that the rest of the book builds on.

\begin{example}{Clocks that run at different rates}\label{ex:lapse}
Suppose we want a region where clocks run at different rates at different
places along $x$, while space itself stays exactly flat. The metric is
\begin{equation}
  \dd s^{2} = -\alpha(x)^{2}\,\dd t^{2} + \dd x^{2} + \dd y^{2} + \dd z^{2},
  \label{eq:lapse-metric}
\end{equation}
where $\alpha(x) > 0$ is the rate of a clock at rest at $x$, relative to
coordinate time. Find the matter this geometry demands.

\solution \emph{Christoffel symbols.} The only non-constant component is
$g_{tt} = -\alpha^{2}$, which depends on $x$. Writing $\alpha' = \dd\alpha/\dd x$,
\cref{eq:christoffel} gives only two independent non-zero symbols:
\begin{equation*}
  \Gamma^{t}{}_{tx} = \Gamma^{t}{}_{xt} = \frac{\alpha'}{\alpha},
  \qquad
  \Gamma^{x}{}_{tt} = \alpha\,\alpha' .
\end{equation*}
\emph{Curvature.} From \cref{eq:riemann}, the non-zero components of the Ricci
tensor are
\begin{equation*}
  R_{tt} = \alpha\,\alpha'', \qquad R_{xx} = -\frac{\alpha''}{\alpha},
\end{equation*}
and the Ricci scalar is $R = g^{tt}R_{tt} + g^{xx}R_{xx} = -2\alpha''/\alpha$.
\emph{Einstein tensor.} From \cref{eq:einstein-tensor},
\begin{align*}
  G_{tt} &= \alpha\alpha'' - \tfrac{1}{2}\left(-\frac{2\alpha''}{\alpha}\right)(-\alpha^{2}) = 0,\\
  G_{xx} &= -\frac{\alpha''}{\alpha} - \tfrac{1}{2}\left(-\frac{2\alpha''}{\alpha}\right) = 0,\\
  G_{yy} &= G_{zz} = -\tfrac{1}{2}\left(-\frac{2\alpha''}{\alpha}\right) = \frac{\alpha''}{\alpha}.
\end{align*}
\emph{Interpretation.} A clock at rest at $x$ has four-velocity
$w^{\mu} = (1/\alpha, 0, 0, 0)$, and by \cref{eq:observed-density} it measures
energy density $\rho = G_{tt}/(8\pi\alpha^{2}) = 0$. The stresses are
\begin{equation}
  \rho = 0, \qquad p_{x} = 0, \qquad p_{y} = p_{z} = \frac{\alpha''}{8\pi\,\alpha}.
  \label{eq:lapse-demand}
\end{equation}
\discussion Changing the rate of clocks from place to place, with space kept
flat, demands \emph{no energy at all}, only stress, and only in the directions
perpendicular to the change. Where $\alpha$ curves upward ($\alpha'' > 0$), the
demand is a pressure; where it curves downward ($\alpha'' < 0$), it is a tension.
If $\alpha$ is a straight line, $\alpha = 1 + gx$, the demand vanishes entirely:
that geometry is flat spacetime described by an observer with constant
acceleration $g$ (\cref{prob:rindler}).
\end{example}

\begin{figure}[htbp]
\centering
\begin{tikzpicture}
\begin{axis}[
  width=0.82\linewidth, height=5.6cm,
  axis lines=middle, xlabel={$x$}, xmin=-3.2, xmax=3.2, ymin=-1.05, ymax=1.75,
  xtick={-2,-1,0,1,2}, ytick={-1,0,1}, legend pos=north east,
  legend style={font=\small, draw=none, fill=white, fill opacity=0.85},
  every axis x label/.style={at={(ticklabel* cs:1.0)}, anchor=west},
]
\addplot[domain=-3.2:3.2, samples=200, very thick, goalblue] {1 + 0.5*exp(-x^2)};
\addlegendentry{clock rate $\alpha(x)$}
\addplot[domain=-3.2:3.2, samples=200, very thick, keygold!80!black]
  {0.5*exp(-x^2)*(4*x^2 - 2)/(1 + 0.5*exp(-x^2))};
\addlegendentry{$\alpha''/\alpha \propto p_{y}$}
\addplot[domain=-0.7071:0.7071, samples=60, fill=red!15, draw=none]
  {0.5*exp(-x^2)*(4*x^2 - 2)/(1 + 0.5*exp(-x^2))} \closedcycle;
\end{axis}
\end{tikzpicture}
\caption{A bump in the rate of clocks, $\alpha = 1 + \tfrac12 e^{-x^{2}}$, and the
transverse stress it demands (\cref{eq:lapse-demand}). Around the peak the
demand is a tension (shaded); on the flanks it is a pressure. The energy
density is zero everywhere.}
\label{fig:lapse-bump}
\end{figure}

\Cref{ex:lapse} raises a fair question. The Earth makes clocks run at different
rates at different heights, yet empty space around the Earth contains no stress.
The difference is that the Earth's geometry also curves space: the factor
$(1 - 2\Phi)$ in front of the spatial part of \cref{eq:weak-field}. In empty
space the curvature of space and the variation of clock rates cancel exactly in
Einstein's equation (\cref{prob:weak-field}). In \cref{ex:lapse} we insisted on
keeping space flat, and the stress is the price of that choice. Designs make
choices like this all the time, and each one has a price that Einstein's
equation computes.

The price is steep. Take a clock-rate change of one part in a billion across one
metre, so that $\alpha'' \sim 10^{-9}\,\si{\metre^{-2}}$. By
\cref{eq:lapse-demand,eq:conversion}, the stress is about
$10^{-9}/(8\pi) \times \SI{1.2e44}{\pascal} \approx \SI{5e33}{\pascal}$. That is
the scale of pressure inside a neutron star, demanded in order to make two clocks
a metre apart differ by one part in a billion.

There is a subtler feature too. Consider light moving along $y$ through the
peak of the bump, where $p_{y} < 0$. In Chapter~5 you will learn that ordinary
matter always has $\rho + p \ge 0$ for the pressure along any direction light
travels. Here $\rho + p_{y} = \alpha''/(8\pi\alpha) < 0$. A simple bump in the rate
of clocks already demands matter that no ordinary material can provide. You
will meet this demand again when we study how each part of a design contributes
to its cost.

\begin{checkpoint}
In \cref{ex:lapse}, suppose instead that clocks run \emph{slowest} at $x = 0$,
with $\alpha = 1 - \tfrac12 e^{-x^{2}}$. Where is the demand a tension, and where
a pressure? Is $\rho + p_{y}$ negative anywhere?
\end{checkpoint}

\begin{example}{A flow of space that shears}\label{ex:shear}
\Cref{ex:flow} described flat spacetime with coordinates that slide at a
constant speed. Now let the speed of the flow vary across space. Consider
\begin{equation}
  \dd s^{2} = -\dd t^{2} + \bigl(\dd x + \beta(y)\,\dd t\bigr)^{2} + \dd y^{2} + \dd z^{2},
  \label{eq:shear-metric}
\end{equation}
a flow of space along $x$ whose speed depends on $y$. Find the energy density
measured by observers who are carried along by the flow.

\solution Observers carried by the flow have $\dd x/\dd t = -\beta$ and feel no
force. Their four-velocity is $w^{\mu} = (1, -\beta, 0, 0)$, which satisfies
$g_{\mu\nu}w^{\mu}w^{\nu} = -1$ (\cref{prob:shear-check}). Carrying out the recipe
for \cref{eq:shear-metric} is longer than for \cref{ex:lapse}, and it is a good
exercise for computer algebra (\cref{prob:shear-check}). The result for the
energy density these observers measure is
\begin{equation}
  \rho = T_{\mu\nu}w^{\mu}w^{\nu} = -\frac{1}{32\pi}\left(\frac{\dd\beta}{\dd y}\right)^{2}.
  \label{eq:shear-density}
\end{equation}
\discussion Where the flow is uniform, $\dd\beta/\dd y = 0$ and nothing is demanded,
just as in \cref{ex:flow}. Wherever the flow \emph{shears}, the observers riding
it measure a \emph{negative} energy density, whichever way the shear goes, since
the derivative is squared. In \cref{ch:what-is} we met warp drives: bubbles of
flat space carried along by a flow of space. At the bubble's wall the flow
changes from the speed of the bubble to zero, so the wall is a region of shear,
and \cref{eq:shear-density} is the reason a warp drive demands negative energy
there. Alcubierre's energy density has exactly this structure
\parencite{Alcubierre1994}.
\end{example}

The numbers are once again enormous. A flow that changes from zero to the speed
of light across one metre has $\dd\beta/\dd y \approx \SI{1}{\metre^{-1}}$, and
\cref{eq:shear-density,eq:conversion} give an energy density of about
$-\SI{1.2e42}{\joule\per\cubic\metre}$, equivalent to a mass density of about
$-\SI{1e25}{\kilogram\per\cubic\metre}$. Nuclear matter, the densest material in
nature, has about $+\SI{2e17}{\kilogram\per\cubic\metre}$.

\section{Every geometry has a price}\label{sec:price}

The recipe of \cref{sec:backwards} never fails. Every metric you can write
down has an Einstein tensor, so every geometry has a demanded stress-energy.
Einstein's equation, read backwards, is a price list: it tells you exactly what
matter any geometry would require.

Whether anything in nature can pay that price is a separate question, and it
is the question the rest of this book is organized around. The examples in this
chapter already show its shape. Shaping the rate of clocks demands stress
without energy, and can demand matter that violates $\rho + p \ge 0$. Shearing
the flow of space demands negative energy density. Chapter~4 splits spacetime
into space and time in the way these examples suggest, with a \emph{lapse} that
sets the rate of clocks and a \emph{shift} that describes the flow of space, and
shows how to compute demands like these efficiently. Chapter~5 states precisely
which demands ordinary matter can meet, and Chapter~6 asks how far quantum
physics can go beyond them.

\begin{chaptersummary}
\item Energy depends on the observer. An observer with four-velocity $w^{\mu}$
  measures the energy density $\rho_{w} = T_{\mu\nu}w^{\mu}w^{\nu}$.
\item The stress-energy tensor $T^{\mu\nu}$ collects energy density, energy
  flux, momentum density and stress into one symmetric array.
\item Dust, perfect fluids, radiation, tension and vacuum energy each have a
  characteristic stress-energy. Tension is negative pressure.
\item Einstein's equation $G_{\mu\nu} = 8\pi T_{\mu\nu}$ (with $G = c = 1$) links
  curvature to matter. Multiply geometric stresses by
  $c^{4}/G = \SI{1.21e44}{\newton}$ to get pascals.
\item Newton's law of gravity is the weak-field limit. Beyond it, pressure
  gravitates through $\rho + 3p$, and tension can make gravity repel.
\item Reading the equation backwards gives the matter any geometry demands.
  Varying clock rates across flat space demands stress without energy; a
  shearing flow of space demands negative energy density.
\end{chaptersummary}

\begin{checkanswers}
\item[3.1] Dust: $\rho + p = \rho$. Radiation: $\rho + p = \tfrac43\rho$. Fluid with
  $p > 0$: $\rho + p > \rho$. Vacuum energy: $\rho + p = 0$. Only vacuum energy
  sits exactly on the boundary.
\item[3.2] Now $\alpha'' > 0$ near $x = 0$, so the demand there is a pressure, and
  on the flanks (where $\alpha$ curves back toward 1 from below) it is a tension.
  Wherever the demand is a tension, $\rho + p_{y} = \alpha''/(8\pi\alpha) < 0$, so
  it is negative on the flanks.
\end{checkanswers}

\begin{problems}
\item \diff{1} A cloud of dust with rest-frame density $\rho_{0}$ moves at speed
  $v$ along $y$. Write down all components of $T^{\mu\nu}$ in the observer's
  frame.
\item \diff{1} \textbf{Geometric units.} Express the mass of the Earth, the Sun
  and a $10^{9}\,\Msun$ black hole in metres. What energy density, in joules per
  cubic metre, corresponds to a geometric energy density of
  \SI{1}{\per\square\kilo\metre}?
\item \diff{2} \textbf{Magnetic tension.} Using the magnetic stress-energy in the
  \emph{Going deeper} box of \cref{sec:catalogue}, verify that
  $\rho + p \ge 0$ for every direction of the pressure, and that the tension along
  the field lines equals the energy density.
\item \diff{2} \textbf{Vacuum energy.} Show from \cref{eq:vacuum-energy} that
  every observer measures the same energy density $\rho_{\text{vac}}$, and that
  the pressure is $-\rho_{\text{vac}}$ in every direction.
\item \diff{2}\label{prob:rho-3p} \textbf{Pressure gravitates.} Einstein's
  equation can be rewritten as $R_{\mu\nu} = 8\pi\bigl(T_{\mu\nu} - \tfrac12 T g_{\mu\nu}\bigr)$,
  with $T = g^{\mu\nu}T_{\mu\nu}$. For a perfect fluid at rest, show that
  $R_{tt} = 4\pi(\rho + 3p)$. Given that $R_{tt} \approx \nabla^{2}\Phi$ in a weak,
  static field, explain why $\rho + 3p$ plays the role of $\rho$ in Poisson's
  equation.
\item \diff{2}\label{prob:weak-field} \textbf{Weak fields.} Compute the Einstein
  tensor of \cref{eq:weak-field} to first order in $\Phi$. Show that
  $G_{tt} \approx 2\nabla^{2}\Phi$, and that all components vanish where
  $\nabla^{2}\Phi = 0$. Explain why the result differs from \cref{ex:lapse}.
\item \diff{2}\label{prob:rindler} \textbf{Uniform acceleration.} Show that the
  metric \cref{eq:lapse-metric} with $\alpha = 1 + gx$ has vanishing Riemann
  tensor. What does an observer at rest at fixed $x$ feel?
\item \diff{2} \textbf{A rippled clock rate.} Take $\alpha = 1 + \epsilon\sin kx$
  with $0 < \epsilon < 1$ in \cref{ex:lapse}. Where is the demand a tension? For
  light moving along $y$, where is $\rho + p_{y}$ negative?
\item \diff{3}\label{prob:shear-check} \textbf{The shearing flow.} Using a computer
  algebra system, compute the Einstein tensor of \cref{eq:shear-metric} and
  verify \cref{eq:shear-density}. Check that $w^{\mu} = (1,-\beta,0,0)$ is a unit
  timelike vector. Then compute the energy density measured by an observer at
  rest in the coordinates, $w^{\mu} = (1,0,0,0)/\sqrt{1-\beta^{2}}$, where
  $|\beta| < 1$, and compare.
\item \diff{3} \textbf{Designing a fast-clock region.} You want a region where
  clocks run twice as fast as clocks far away, with space kept flat and the
  transition taking place over \SI{10}{\metre}. Using \cref{ex:lapse}, sketch a
  suitable $\alpha(x)$, locate the tensions and pressures it demands, and estimate
  their size in pascals.
\end{problems}

\begin{furtherreading}
\item[\textcite{Schutz2009}] A careful introduction to the stress-energy tensor
  through dust and perfect fluids.
\item[\textcite{Carroll2004}] Einstein's equation, its Newtonian limit and the
  Bianchi identity at graduate level.
\item[\textcite{Hartle2003}] A physical approach to Einstein's equation and to
  what curvature does.
\item[\textcite{Morris1988a}] The engineer's calculation of this chapter,
  applied to traversable wormholes.
\item[\textcite{Alcubierre1994}] The original warp drive. After
  \cref{ex:shear} you can follow its energy density calculation.
\end{furtherreading}
